\section{The Salvetti complex and spherical parabolic cosets}
\label{sec:topology}

We first specify the space whose universal cover will be contracted.
Let $A=A(W,S)$ be the Artin group in Theorem~\ref{thm:main}, and let
\[
  \vartheta\colon A\longrightarrow W,
  \qquad \vartheta(\sigma_s)=s,
\]
be its Coxeter quotient.  For $T\subseteq S$, write $W_T=\langle T\rangle$
and $A_T=\langle\sigma_s:s\in T\rangle$, and put
\[
  \mathcal S^{\mathrm f}
  =\{T\subseteq S:W_T\text{ is finite}\}.
\]
The members of $\mathcal S^{\mathrm f}$ are called \emph{spherical}.
In particular, $\varnothing\in\mathcal S^{\mathrm f}$ and
$A_{\varnothing}=\{1\}$.

For a poset $\mathcal P$, its realization $|\mathcal P|$ will always mean
the order complex: vertices are elements of $\mathcal P$, simplices are
nonempty finite chains, and the realization has its CW topology.

\subsection{Classical inputs}

We use the following facts about Coxeter systems
\cite[Chapter IV, \S1, nos.~4--5 and 8]{Bourbaki2002}.
Each $(W_T,T)$ is the Coxeter
system defined by the submatrix on $T$.  A left coset of $W_R$ in
$W_T$, for $R\subseteq T$, has a unique element $u_0$ of minimal length,
and
\begin{equation}
  \ell(u_0v)=\ell(u_0)+\ell(v)
  \qquad(v\in W_R).
  \label{eq:topology-length-additivity}
\end{equation}
Reduced expressions for an element of $W$ are related by braid moves.
Consequently, if $w=s_1\cdots s_r$ is reduced, the element
\begin{equation}
  \iota(w)=\sigma_{s_1}\cdots\sigma_{s_r}\in A
  \label{eq:topology-positive-lift}
\end{equation}
is independent of the reduced expression.  Whenever
$\ell(uv)=\ell(u)+\ell(v)$, it satisfies
$\iota(uv)=\iota(u)\iota(v)$.

We also use the geometric representation and its chamber description:
the standard geometric form is positive definite exactly when the
Coxeter group is finite
\cite[Chapter V, \S4, no.~8, Theorem 2]{Bourbaki2002};
in finite type the chambers form a simply
transitive Coxeter orbit, signs of roots detect descents, and the
translates of a closed fundamental chamber cover the dual space
\cite[Chapter V, \S3, no.~3, and \S4, nos.~4 and 6]{Bourbaki2002}.
These assertions apply both to spherical submatrices of the original
system and to the finite Coxeter systems constructed below.

\begin{theorem}[Standard parabolic embedding]
\label{thm:topology-parabolic-embedding}
For every $T\subseteq S$, the homomorphism from the Artin group defined
by the submatrix on $T$ to $A$, sending its standard generators to the
corresponding $\sigma_t$, is injective.  Its image is $A_T$.
\end{theorem}

This is van der Lek's standard parabolic theorem
\cite[Chapter II, Section 4, Theorem 4.13(i)]{VanDerLek1983}.
No spherical hypothesis is imposed here.

The positive monoid on a subset $T$ means the monoid with the same
generators and positive braid relations as its Artin group.  Its
identity is the empty word.  An element $b$ \emph{right-divides} $a$
if $a=cb$ for some positive element $c$; left divisibility is defined
in the opposite order.

\begin{theorem}[Finite-type positive monoids]
\label{thm:topology-finite-monoid}
Suppose that $T$ is spherical, and let $A_T^+$ be its positive Artin
monoid.  Then the following statements hold.
\begin{enumerate}[label=\textup{(\roman*)}]
  \item The natural map $A_T^+\to A_T$ is injective.  In particular,
    $A_T^+$ is cancellative on both sides.
  \item Every element of $A_T$ can be written both as $ab^{-1}$ and
    as $c^{-1}d$, with $a,b,c,d\in A_T^+$.
  \item Any two elements of $A_T^+$ have a greatest common right
    divisor, and also a greatest common left divisor.
  \item If $w_T$ is the longest element of $W_T$, then
    $\Delta_T=\iota(w_T)$ is the finite-type fundamental element.
    Its left divisors and its right divisors coincide and are exactly
    the elements $\iota(w)$, $w\in W_T$.
\end{enumerate}
\end{theorem}

These are the finite-type monoid results of Brieskorn--Saito and
Deligne.  Cancellation and fractions are given in
\cite[Propositions 2.3 and 5.5(iii)--(iv), and Theorem 5.6]{BrieskornSaito1972} and
\cite[Theorem 4.14]{Deligne1972}; the gcd statement is
\cite[Proposition 4.2]{BrieskornSaito1972}.  For the fundamental
element and its divisors, see
\cite[Lemma 5.4 and Propositions 5.5 and 5.7]{BrieskornSaito1972}.
We use Theorem~\ref{thm:topology-finite-monoid} only for spherical
subsets and for finite Coxeter unfoldings.  In particular, fractions
in a positive monoid are never assumed for the full system when it
is not spherical.

\subsection{Permutahedral cells and their face maps}

For each nonempty spherical $T$, take a point in the interior of a
fundamental chamber of the finite reflection representation of $W_T$.
Choose these points compatibly with diagram components: on a reducible
system use the direct sum of the points chosen for its components.
The convex hull of its $W_T$-orbit is a Coxeter permutahedron $\Pi_T$.
It is a closed ball of dimension $|T|$, with vertices labeled by
$W_T$ and with a distinguished \emph{source} labeled $1$.
Set $\Pi_{\varnothing}$ to be a point.

The nonempty faces of $\Pi_T$ are indexed by cosets $u_0W_R$, where
$R\subseteq T$ and $u_0$ is the shortest element of its coset.
The face with this index will be denoted $F_{u_0,R}$; its vertices
have labels $u_0v$, $v\in W_R$.  There is a specified cellular
homeomorphism
\begin{equation}
  j_{u_0,R}\colon \Pi_R\longrightarrow F_{u_0,R},
  \qquad v\longmapsto u_0v
  \quad\text{on vertices}.
  \label{eq:topology-face-identification}
\end{equation}
One precise choice is obtained by taking the barycentric subdivisions
of both face lattices, sending the center of the face $v_0W_Q$ to the
center of $u_0v_0W_Q$, and extending affinely over their simplices.
Centers on product faces are chosen as products of the centers of
their factors.

These identifications are compatible on nested faces.  Indeed, if
$v_0$ is the shortest representative of a left $W_Q$-coset in $W_R$,
then \eqref{eq:topology-length-additivity} makes $u_0v_0$ the shortest
representative of the corresponding left $W_Q$-coset in $W_T$.
On barycentric subdivisions this gives the equality
\begin{equation}
  j_{u_0,R}\circ j_{v_0,Q}=j_{u_0v_0,Q}.
  \label{eq:topology-nested-faces}
\end{equation}

Define $X(W,S)$ by taking one copy of $\Pi_T$ for every
$T\in\mathcal S^{\mathrm f}$ and, for every proper face
$F_{u_0,R}\subset\Pi_T$, identifying $j_{u_0,R}(y)$ with the point
$y$ in the copy of $\Pi_R$.  Equation~\eqref{eq:topology-nested-faces}
makes all these identifications compatible.  The resulting CW complex
has one open $|T|$-cell of each spherical type $T$.

An edge in a model cell whose endpoints are $u_0$ and $u_0s$, with
$\ell(u_0s)=\ell(u_0)+1$, projects to the edge of type $s$, oriented
positively.  This convention fixes the attaching maps, including their
orientations.  Equivalently, the Coxeter Salvetti complex has cells
$(w,T)$, $w\in W$, $T\in\mathcal S^{\mathrm f}$, with the same models
and face parameters $(wu_0,R)$.  Its $W$-action is left multiplication
on $w$, and its quotient has precisely the identifications just given.
Thus this is the oriented Coxeter--Salvetti model
\cite[Theorem~1.4]{Salvetti1994}; compare
\cite[Theorem~3.3 and Section~5]{Paolini2026}.

The one-skeleton has one vertex and one oriented edge $\sigma_s$ for
each $s\in S$.  If $\{s,t\}$ is spherical, its model is a $2m_{st}$-gon.
The two increasing paths from its source to its opposite vertex read
the two alternating words of length $m_{st}$.  Its attaching relation
is their equality.  If $m_{st}=\infty$, there is no cell of this type.
It follows that the two-skeleton is the Artin presentation complex,
and hence
\begin{equation}
  \pi_1(X(W,S))=A.
  \label{eq:topology-fundamental-group}
\end{equation}
The cells of dimension at least three remain part of $X(W,S)$.

We invoke the spherical theorem in the following precise form.

\begin{theorem}[Spherical asphericity]
\label{thm:topology-spherical-asphericity}
If $T$ is spherical, the universal cover of the CW complex
$X(W_T,T)$ with the preceding permutahedral attaching maps is
contractible.
\end{theorem}

Deligne proved asphericity of the corresponding finite reflection
arrangement complement \cite[opening theorem, p.~273]{Deligne1972}.
The Salvetti formulation is given, for example, in
\cite[Theorem 3.3 and Corollary 5.7]{Paolini2026}.
The hypothesis here is finiteness of $W_T$, not merely finiteness
of its generating set.

\subsection{The regular universal cover}

Let $\widetilde X$ be the universal cover of $X(W,S)$, and label its
vertices by $A$, with the vertex $1$ chosen over the unique vertex of
$X(W,S)$.  A positive edge of type $s$ goes from $\alpha$ to
$\alpha\sigma_s$.  A lift of a cell of type $T$ is uniquely determined
by the lift $\alpha$ of its source, so denote this open cell by
$(\alpha,T)$.

A reduced gallery from the source of $\Pi_T$ to its vertex $u$ reads
a reduced word for $u$.  Its lift therefore ends at $\alpha\iota(u)$.
In particular, the face indexed by $u_0W_R$ has parameter
\begin{equation}
  (\alpha\iota(u_0),R),
  \qquad R\subseteq T,
  \quad u_0\text{ shortest in }u_0W_R\subseteq W_T.
  \label{eq:topology-lifted-faces}
\end{equation}
The corresponding characteristic map on that face is
\eqref{eq:topology-face-identification}.  Compatibility also follows
directly from
$\iota(u_0v_0)=\iota(u_0)\iota(v_0)$.

The lifted CW structure is regular.  To see this, consider two open
faces of a single $\Pi_T$.  If their types differ, their lifted cells
differ.  If both have type $R$, with shortest representatives $u_0$
and $u_1$, equality of their lifted parameters would give
$\alpha\iota(u_0)=\alpha\iota(u_1)$.  Applying $\vartheta$ gives
$u_0=u_1$, so the faces coincide.  Each open face maps
homeomorphically to its lifted cell.  It follows inductively on
dimension that the characteristic map of $\Pi_T$ is injective on its
whole closed ball.  Since the model is compact and the cover is
Hausdorff, this map is an embedding.  Thus the face poset
$\mathcal F(\widetilde X)$ has realization canonically homeomorphic
to the barycentric subdivision of $\widetilde X$.

\subsection{The spherical-coset comparison}

The modified Deligne complex of Charney--Davis is built from spherical
parabolic cosets~\cite[Section~1.5]{CharneyDavis1995}.  We use its
right-coset presentation.  Define
\begin{equation}
  \mathcal D
  =\{A_Tx:x\in A,\ T\in\mathcal S^{\mathrm f}\},
  \label{eq:topology-deligne-poset}
\end{equation}
with types retained, ordered by coset inclusion with nested types.
Thus $A_Ry\le A_Tx$ means $R\subseteq T$ and
$A_Ry\subseteq A_Tx$.  Inversion identifies this poset with the
left-coset poset
$\mathcal D^{\ell}=\{xA_T:x\in A,\ T\in\mathcal S^{\mathrm f}\}$.

We use the following form of Quillen's Theorem A
\cite[Section 1, Theorem A and its dual formulation]{Quillen1973}:
if $f\colon\mathcal P\to\mathcal Q$ is a map of small posets and
\[
  \bigl|f^{-1}(\mathcal Q_{\le q})\bigr|
  \quad\text{is contractible for every }q\in\mathcal Q,
\]
then $|f|\colon|\mathcal P|\to|\mathcal Q|$ is a homotopy equivalence.
Here $\mathcal Q_{\le q}=\{q'\in\mathcal Q:q'\le q\}$ is the whole
lower ideal, and the posets and fibers need not be finite.

\begin{lemma}[Salvetti--coset comparison]
\label{lem:salvetti-comparison}
There is a homotopy equivalence
\[
  \widetilde X(W,S)\simeq|\mathcal D|.
\]
Consequently, contractibility of $|\mathcal D|$ implies the conclusion
of Theorem~\ref{thm:main}.
\end{lemma}

\begin{proof}
Define
\[
  f\colon\mathcal F(\widetilde X)\longrightarrow\mathcal D^{\ell},
  \qquad f(\alpha,T)=\alpha A_T.
\]
Formula~\eqref{eq:topology-lifted-faces} makes $f$ order preserving:
if $(\alpha\iota(u_0),R)$ is a face of $(\alpha,T)$, then
$\iota(u_0)\in A_T$ and $R\subseteq T$, so
$\alpha\iota(u_0)A_R\subseteq\alpha A_T$.

Fix $\beta A_T\in\mathcal D^{\ell}$.  The inverse image of its
lower ideal consists exactly of the cells
\begin{equation}
  (\alpha,R),\qquad \alpha\in\beta A_T,\quad R\subseteq T.
  \label{eq:topology-fiber-cells}
\end{equation}
Indeed, the indicated conditions imply
$\alpha A_R\subseteq\beta A_T$, and the converse follows from the
retained types and from $\alpha\in\alpha A_R$.

By Theorem~\ref{thm:topology-parabolic-embedding}, write
$\alpha=\beta h$ uniquely with $h$ an element of the Artin group on
$T$.  Under this correspondence, every face formula becomes
\[
  (h,R)\longmapsto(h\iota(u_0),Q),
  \qquad Q\subseteq R\subseteq T,
\]
with $u_0$ the shortest representative of the specified $W_Q$-coset
in $W_R$.  These are exactly the lifted face maps of $X(W_T,T)$.
Thus the poset in \eqref{eq:topology-fiber-cells} is its universal-cover
face poset.  Its realization is contractible by
Theorem~\ref{thm:topology-spherical-asphericity}.  For $T=\varnothing$
it is a single vertex.

Quillen's Theorem A now gives
$|\mathcal F(\widetilde X)|\simeq|\mathcal D^{\ell}|$.
Regularity identifies the left side with a subdivision of
$\widetilde X$, and inversion identifies the right side with
$|\mathcal D|$.  This proves the lemma.
\end{proof}

\begin{remark}
The comparison uses the full spherical universal covers over lower
ideals, including all their higher cells.  It does not assert that the
translation action on $\mathcal D$ is free: its vertex stabilizers are
conjugates of standard parabolic subgroups.  The classical comparison appears in
\cite[Corollary~3.2.6]{CharneyDavis1995}; it is also established in the
more general framework of
\cite[Section 3, Theorem 3.1 and Lemma 3.2]{GodelleParis2012}.
\end{remark}

We record two elementary coset facts for the later link calculations.
If $R\subseteq T$ and $A_Ry$ meets $A_Tx$, then
$A_Ry\subseteq A_Tx$.  For a fixed coset $A_Ry$ and a specified
larger spherical type $T$, its unique upper extension is $A_Ty$.
Thus lower intervals can be computed entirely inside the indicated
parabolic subgroup, and upper extensions are determined by their
types.  These assertions use only inclusions with nested types.

Finally, if $T=T_1\sqcup\cdots\sqcup T_r$ is the decomposition into
components of its Coxeter diagram, the presentation and
Theorem~\ref{thm:topology-parabolic-embedding} give
\begin{equation}
  A_T=A_{T_1}\times\cdots\times A_{T_r}.
  \label{eq:topology-parabolic-product}
\end{equation}
Cosets of subgroups whose types lie in these blocks, and their
inclusion relations, are accordingly products.  When $S$ is empty,
both $X(W,S)$ and $|\mathcal D|$ are points; in the rest of the proof
we may assume $S\ne\varnothing$.
